% Part: incompleteness
% Chapter: incompleteness-provability
% Section: provability-conditions

\documentclass[../../../include/open-logic-section]{subfiles}

\begin{document}

% SEGMENT OLP-0318-S01

\olfileid{inc}{inp}{prc}

\olsection{$\Th{PA}$ க்கான \usetoken{S}{derivability} நிபந்தனைகள்}

Peano எண்கணிதம், அல்லது $\Th{PA}$, என்பது எல்லா !!{formula}s க்குமான
தூண்டல் அடிகோள்களைச் சேர்த்து $\Th{Q}$ ஐ நீட்டிக்கும் கோட்பாடாகும்.
வேறு விதமாகச் சொன்னால், பின்வரும் வடிவிலான அடிகோள்களை ஒவ்வொரு
!!{formula}~$!A$ க்கும் $\Th{Q}$ உடன் சேர்க்கிறோம்:
\[
(!A(0) \land \lforall[x][(!A(x) \lif !A(x'))]) \lif \lforall[x][!A(x)]
\]
இது உண்மையில் ஓர் \emph{திட்டம்} என்பதை நினைவில் கொள்க; அதாவது
எண்ணற்ற அடிகோள்கள் இதில் உள்ளன (மேலும் $\Th{PA}$ வரையறுக்கப்பட்ட
எண்ணிக்கையால் அடிகோளாக்கத்தக்கது {\em அல்ல} என்று பின்னர் தெரியும்).
ஆயினும், ஒரு குறியீடுகளின் சரம் தூண்டல் அடிகோளின் நிகழ்வா இல்லையா
என்பதைச் செயல்வழியாகத் தீர்மானிக்க முடியும். ஆகவே $\Th{PA}$ இன்
அடிகோள்களின் கணம் கணிக்கத்தக்கது. $\Th{PA}$ ஆனது $\Th{Q}$ விட மிகவும்
வலுவான கோட்பாடாகும். எடுத்துக்காட்டாக, வழக்கமான முறையில் தூண்டலைப்
பயன்படுத்தி கூட்டலும் பெருக்கலும் பரிமாற்றுச் செயலிகள் என்பதை எளிதாக
நிறுவலாம். உண்மையில், பெரும்பாலான இறுதிக்குட்பட்ட எண்ணியல் மற்றும்
கணிதக்கூற்றுக் காரணங்கூறல்களை $\Th{PA}$ இல் மேற்கொள்ள முடியும்.

$\Th{PA}$ கணிக்கத்தக்கவாறு அடிகோளாக்கப்பட்டதால், அதன்
!!{derivability} பயனிலை $\Prf[\Th{PA}](x,y)$ கணிக்கத்தக்கது; எனவே அது
$\Th{Q}$ இல் (அதனால் $\Th{PA}$ இலும்) பிரதிநிதித்துவப்படுத்தப்படுகிறது.
முன்பைப் போலவே, அந்த உறவைப் பிரதிநிதித்துவப்படுத்தும் வாய்பாட்டைக்
குறிக்க $\OPrf[\Th{PA}](x,y)$ என்று எழுதுவோம். மேலும்
$\OProv[\Th{PA}](y)$ என்பது
$\lexists[x][\Prf[\Th{PA}](x,y)]$ என்ற வாய்பாடாக இருக்கட்டும்.
உள்ளுணர்வாக, இது “$y$ என்பது $\Th{PA}$ இன் அடிகோள்களிலிருந்து
!!{derivable}” என்று கூறுகிறது. $\Th{Q}$ இன் அடிகோள்களைவிடச் சிறிது
கூடுதல் நமக்கு தேவைப்படுவதன் காரணம், நாம் பயன்படுத்தும் கோட்பாடு இந்த
!!{derivability} பயனிலை பற்றிய சில அடிப்படை உண்மைகளை !!{derive}
செய்யும் அளவுக்கு வலுவாக இருக்க வேண்டும் என்பதே. தேவையான உண்மைகள்
பின்வருமாறு:
\begin{enumerate}
\item[P1.] $\Th{PA} \Proves !A$ எனில்,
  $\Th{PA} \Proves \OProv[\Th{PA}](\gn{!A})$.
\item[P2.] எல்லா !!{formula}s~$!A$ மற்றும்~$!B$ க்கும்,
  \[
  \Th{PA} \Proves \OProv[\Th{PA}](\gn{!A \lif !B}) \lif
  (\OProv[\Th{PA}](\gn{!A}) \lif \OProv[\Th{PA}](\gn{!B})).
  \]
\item[P3.] ஒவ்வொரு !!{formula}~$!A$ க்கும்,
  \[
  \Th{PA} \Proves \OProv[\Th{PA}](\gn{!A})
  \lif \OProv[\Th{PA}](\gn{\OProv[\Th{PA}](\gn{!A})}).
  \]
\end{enumerate}
இந்த மூன்று பண்புகள் பொருந்துகின்றன என்பதைச் சரிபார்க்க ஒரே வழி,
$\OProv[\Th{PA}](y)$ என்ற !!{formula} வைத் துல்லியமாக விவரித்து,
தொடர்புடைய முறைசார் !!{derivation}s ஐ விவரிக்க $\Th{PA}$ இன்
அடிகோள்களைப் பயன்படுத்துவதாகும். (1), (2) ஆகிய நிபந்தனைகள் எளிதானவை;
உண்மையான வேலை (3) ஆம் நிபந்தனைக்குத்தான் தேவைப்படுகிறது. (அது
எத்தகைய வேலையை உட்படுத்தும் என்று சிந்தித்துப் பாருங்கள்~\dots)
விவரங்களை முழுமையாக எழுதுவது சலிப்பூட்டுவதும் பயனற்றதுமாக இருக்கும்;
ஆகவே $\Th{PA}$ மேற்கண்ட மூன்று நிபந்தனைகளையும் கொண்டுள்ளது என்று
நம்பிக்கையுடன் ஏற்றுக்கொள்ளுமாறு இங்கு கேட்கிறோம். நியாயமான முறையில்
தேர்ந்தெடுக்கப்பட்ட $\OProv[\Th{PA}](y)$ பின்வரும் நிபந்தனையையும்
நிறைவேற்றும்:
\begin{enumerate}
\item[P4.] $\Th{PA} \Proves \OProv[\Th{PA}](\gn{!A})$ எனில்,
  $\Th{PA} \Proves !A$.
\end{enumerate}
ஆனால் இந்த உண்மை நமக்கு தேவைப்படாது.

\begin{digress}
தற்செயலாக, நாமும் இப்போது செய்வதைப் போலவே கோடலும் சோம்பேறியாக
இருந்தார். 1931 ஆம் ஆண்டுக் கட்டுரையின் முடிவில் அவர் இரண்டாம்
முழுமையின்மைத் தேற்றத்தின் நிறுவலைச் சுருக்கமாகக் காட்டி, விவரங்களை
பின்னர் வரும் கட்டுரையில் தருவதாக உறுதியளித்தார். அதை எழுத அவருக்கு
நேரம் கிடைக்கவில்லை; வாதத்தைப் புரிந்த அனைவரும் அதை முழுமையாகச்
செய்ய முடியும் என்று நம்பியதால், விவரங்களை நிரப்ப வேண்டிய அவசியம்
அவருக்குத் தோன்றவில்லை.
\end{digress}

\end{document}
